% Generated draft: requires mathematical and source review.
\documentclass[11pt]{article}
\usepackage[margin=25mm]{geometry}
\usepackage{fontspec}
\setmainfont{Latin Modern Roman}
\usepackage{amsmath,amssymb,mathtools,unicode-math}
\setmathfont{Latin Modern Math}
\usepackage{enumitem,needspace,xurl,hyperref}
\setlength{\parindent}{0pt}
\setlength{\parskip}{5pt}
\setlength{\emergencystretch}{3em}
\newcommand{\N}{\mathbb{N}}
\newcommand{\Z}{\mathbb{Z}}
\newcommand{\Q}{\mathbb{Q}}
\newcommand{\R}{\mathbb{R}}
\newcommand{\C}{\mathbb{C}}
\newcommand{\F}{\mathbb{F}}
\newcommand{\A}{\mathbb{A}}
\newcommand{\PP}{\mathbb{P}}
\newcommand{\E}{\mathbb{E}}
\newcommand{\eps}{\varepsilon}
\newcommand{\dd}{\,\mathrm{d}}
\begin{document}
\section*{UnsolvedMath ID: 30003099}
\textbf{Source problem code:} catalogue entry
\subsection{Definitions and mathematical statement}
The supplied record asks for a valuation-dependent construction in situations where embedded resolution for a variety \(X\) is not known to exist. This describes a limitation of available knowledge, not an assumption that embedded resolution is impossible. The ground field, its characteristic, the dimension and integrality hypotheses on \(X\), and the class of valuations are not specified. Consequently, the following describes the intended construction and its conditions without asserting a universally quantified theorem under invented hypotheses.

For clarity, in an instance where \(X\) is an integral variety over a field \(k\), write \(K=k(X)\). A valuation of \(K\) is a map \(v:K^{\times}\to\Gamma\), where \(\Gamma\) is a totally ordered abelian group, satisfying
\[
v(ab)=v(a)+v(b),\qquad v(a+b)\geq\min\{v(a),v(b)\}\quad(a+b\ne0).
\]
Its valuation ring and maximal ideal are
\[
V_v=\{0\}\cup\{a\in K^{\times}:v(a)\geq0\},\qquad
\mathfrak m_v=\{0\}\cup\{a\in K^{\times}:v(a)>0\}.
\]
A center of \(v\) on an integral model \(Y\) with function field \(K\) is a point \(y\in Y\) such that
\[
\mathcal O_{Y,y}\subseteq V_v,\qquad
\mathfrak m_y=\mathfrak m_v\cap\mathcal O_{Y,y}.
\]
The source does not specify whether the chosen valuation is initially centered on \(X\), whether it is trivial on \(k\), or what restrictions are imposed on its rank or residue field.

A toric variety \(Z\) has a dense algebraic torus \(T\) whose action extends to \(Z\); its toric boundary is \(D=Z\setminus T\), with reduced structure. A toric birational map is a birational map compatible with the torus actions. In a presentation by a birational morphism \(\pi:Z'\to Z\), let \(U\subseteq Z\) be an open set over which \(\pi\) is an isomorphism. If a locally re-embedded model \(X_0\subseteq Z\) meets \(U\) densely, its strict transform \(X'\subseteq Z'\) is the closure of \(\pi^{-1}(X_0\cap U)\), taken with reduced structure. The source does not specify whether its maps must be proper morphisms or which local re-embeddings are allowed.

\textbf{Requested construction.} Given the variety \(X\) and a chosen valuation in the intended, unspecified class, find suitable local re-embeddings and toric birational maps such that the resulting strict transform \(X'\), birationally identified with \(X\), has a center \(x'\) of the chosen valuation, is nonsingular at \(x'\), and is transverse there to the toric boundary of the final ambient toric variety.

The regularity condition can be expressed as
\[
\mathcal O_{X',x'}\text{ is a regular local ring}.
\]
The boundary condition needs an additional convention. For example, when \(X'\) and the ambient \(Z'\) are smooth near the relevant point, transversality to a smooth boundary stratum \(S\) through a geometric point \(p\) above \(x'\) means
\[
T_pX'+T_pS=T_pZ'.
\]
Requiring this for all boundary strata through the point is one possible precise interpretation. The supplied record does not specify this convention, ambient smoothness, a normal crossings requirement, or whether nonsingularity is intended to mean regularity or smoothness over the ground field. These qualifications remain unresolved; the displayed interpretation is not asserted to be the exact original formulation.
\subsection{Short English statement}
For a chosen valuation of a variety, seek local re-embeddings and toric birational modifications whose strict transform is nonsingular and meets the toric boundary transversely at the valuation's center. The intended fields, varieties, valuations, allowed modifications, and meaning of boundary transversality are unspecified. The supplied report assessment presents this as a possible approach where embedded resolution is unavailable, rather than a general solution.
\subsection{Sources}
\begin{itemize}
\item Supplied local UnsolvedMath registry record: permanent ID \(30003099\), source problem code \texttt{catalogue entry}, title \emph{Toric Local Uniformization Without Embedded Resolution}, in the collection \emph{Oberwolfach Reports Open Problems}. The record gives proposed year \(2016\) and names no proposer.
\item Supplied report citation: \emph{Toric Geometry (2017)}, Oberwolfach Reports, DOI \texttt{10.4171/owr/2016/19}, pp.\ 963--1026. Supplied source URL: \url{https://doi.org/10.4171/owr/2016/19}. The background describes Oberwolfach Reports as documenting mathematical workshops and their open-problem discussions. No precise page or individual contribution is identified for this problem.
\item Supplied literature source, labeled ``OWR 2016 report'': \url{https://doi.org/10.4171/owr/2016/19}. The literature assessment was checked on \texttt{2026-08-21}. This is supplied dated curation, not independently verified here. It labels the problem open and describes valuation-based toric uniformization as a possible approach when embedded resolution is unavailable, not as a general solution. The supplied background expressly calls for independent verification of time-sensitive status and cited claims.
\item The registry labels the statement \texttt{corrected\_verified}, explaining that mathematical content was preserved while notation, typography, or omitted definitions were repaired for self-contained publication. This is supplied editorial metadata, not an independent mathematical verification of this draft or of the problem's status.
\end{itemize}
\end{document}
